\documentclass{article}
\usepackage[english,russian]{babel}
\usepackage{textcomp}
\usepackage{geometry}
  \geometry{left=2cm}
  \geometry{right=1.5cm}
  \geometry{top=1.5cm}
  \geometry{bottom=2cm}
\usepackage{tikz}
\usepackage{multicol}
\usepackage{hyperref}
\usepackage{listings}
\pagenumbering{gobble}

\lstdefinestyle{csMiptCppStyle}{
  language=C++,
  basicstyle=\linespread{1.1}\ttfamily,
  columns=fixed,
  fontadjust=true,
  basewidth=0.5em,
  keywordstyle=\color{blue}\bfseries,
  commentstyle=\color{gray},
  texcl=true,
  stringstyle=\ttfamily\color{orange!50!black},
  showstringspaces=false,
  numbersep=5pt,
  numberstyle=\tiny\color{black},
  numberfirstline=true,
  stepnumber=1,      
  numbersep=10pt,
  backgroundcolor=\color{white},
  showstringspaces=false,
  captionpos=b,
  breaklines=true
  breakatwhitespace=true,
  xleftmargin=.2in,
  extendedchars=\true,
  keepspaces = true,
  tabsize=4,
  upquote=true,
}


\lstdefinestyle{csMiptCppLinesStyle}{
  style=csMiptCppStyle,
  frame=lines,
}

\lstdefinestyle{csMiptCppBorderStyle}{
  style=csMiptCppStyle,
  framexleftmargin=5mm, 
  frame=shadowbox, 
  rulesepcolor=\color{gray}
}


\lstdefinestyle{csMiptBash}{
breaklines=true,
frame=tb,
language=bash,
breakatwhitespace=true,
alsoletter={*()"'0123456789.},
alsoother={\{\=\}},
basicstyle={\ttfamily},
keywordstyle={\bfseries},
literate={{=}{{{=}}}1},
prebreak={\textbackslash},
sensitive=true,
stepnumber=1,
tabsize=4,
morekeywords={echo, function},
otherkeywords={-, \{, \}},
literate={\$\{}{{{{\bfseries{}\$\{}}}}2,
upquote=true,
frame=none
}




\lstset{style=csMiptCppLinesStyle}
\lstset{literate={~}{{\raisebox{0.5ex}{\texttildelow}}}{1}}


\renewcommand{\thesection}{\arabic{section}}
\makeatletter
\def\@seccntformat#1{\@ifundefined{#1@cntformat}%
   {\csname the#1\endcsname\quad}%    default
   {\csname #1@cntformat\endcsname}}% enable individual control
\newcommand\section@cntformat{Часть \thesection:\space}
\makeatother




\begin{document}
\title{Семинар \#2: CMake \vspace{-5ex}}\date{}\maketitle

\section{Системы сборки}
\begin{description}
\item \textbf{Система сборки} (англ. \textit{build system}) -- это программа (или набор программ), предназначенная для автоматизации преобразования исходного кода в исполняемые файлы или библиотеки. 
\end{description}
Для языка C++ существуют множество систем сборки, вот самые популярные из них:


\begin{itemize}
\item \textbf{CMake} -- используется в абсолютном большинстве проектов на C++.\\ Является не просто системой сборки, но и генератором для других систем.
\item \textbf{Make} -- классическая система сборки, часто используется в ОС Linux для языка C.
\item \textbf{Meson} -- альтернатива CMake, но менее популярна.
\item \textbf{Bazel} -- система сборки для разных языков (C++, Java, Python), разработанная Google.
\item \textbf{Ninja} -- минималистичная система сборки, ориентированная на скорость.
\item \textbf{MSBuild} -- система сборки для Windows, используется в IDE Visual Studio.
\end{itemize}



\subsection*{Задачи, которые решают системы сборки}
Представьте, что есть проект, написанный на C++, который содержит сотни файлов и зависит от многих сторонних библиотек. Как его скомпилировать? Можно было бы просто вызвать компилятор \texttt{g++} из командной строки и передать ему все файлы проекта, информацию о всех подключаемых библиотеках и все опции компиляции и линковки необходимые для проекта. В теории так можно скомпилировать даже большой проект, однако такой подход имеет множество недостатков. Например, вводимая команда будет очень большой и в ней можно будет легко  запутаться. Но это только одна из многих проблем данного подхода. 

\begin{itemize}
\item \textbf{Организация структуры проекта}

\item \textbf{Множество исполняемых файлов/библиотек в одном проекте}

\item \textbf{Инкрементальная сборка}

\item \textbf{Управление зависимостями}

\item \textbf{Кроссплатформенность}

\item \textbf{Поддержка различных компиляторов}

\item \textbf{Поддержка различных IDE}


\item \textbf{Настройка опций компиляции и линковки}

\item \textbf{Генерация дополнительных файлов проекта}

\item \textbf{Тестирования и установка}
\end{itemize}


\subsection*{Что такое CMake?}
\begin{itemize}
\item CMake -- это генератор систем сборки.
\item CMake -- это система сборки.
\item CMake -- это язык программирования, для описания сборки проекта.
\end{itemize}


\subsection*{Сборка HelloWorld с помощью CMake}

\subsection*{CMake как генератор систем сборки}


\newpage
\section{Простые проекты с использованием CMake}

\section{Таргеты}


\section{CMake как язык программирования}


\end{document}
